% --- 背景 ---
\section{背景}
本文を読み進めるうえで必要となる事前知識や，この研究の先行研究について説明する．
この研究に強く関わるものはより詳しく説明する．
\begin{itemize}
    \item 単なる知識の羅列（教科書の写し）にしない．提案手法で使う数式や定義を厳選して定義する．使わない定義は書かない．
    \item 確率変数$X$は大文字，実現値$x$は小文字，ベクトル$\bm{x}$は太字など，記号のルールをここで確定させ，論文を通して一貫する．
    \item 既存手法の限界を突く．「手法Aは画期的だが，〇〇の条件下では動かない」→「だから私の研究が必要」という流れを作る．
\end{itemize}

\subsection{小節}
小節2.1を書く．

\subsubsection{小々節}
小節2.1.1を書く．こまめに節や小節に分け，見やすさを意識しよう．